% =====================================================================================================
\section{The proved results}\label{sec:proved}\label{sec:results}

Six results are proved. The first two are about the posterior the software targets; the next three are about what a run
of the software reports; the last says that the rate of the first cannot be improved. Each is given in the same
format: \emph{in words}; the formal statement, which cites the assumptions of Section~\ref{sec:assumptions} by name; a
\emph{status} box, outside the statement, quoting the project's registry; one paragraph on \emph{why it is natural}; and
a one-line \emph{scope} (which of our designs are inside). Table~\ref{tab:results} summarises them.
Section~\ref{sec:planned} lists the results that are planned and not yet proved.

\newcommand{\whynat}[1]{\par\noindent\textit{Why it is natural.} #1\par\smallskip}
\newcommand{\vssl}[1]{\par\noindent\textit{Compared with \SL.} #1\par\smallskip}

\begin{table}[!t]
\centering\footnotesize
\setlength{\tabcolsep}{3.5pt}
\renewcommand{\arraystretch}{1.12}
\caption{The six proved results. Statuses are those of the project's registry (Section~\ref{sec:statuses}); every check
is independent of the author.}\label{tab:results}
\begin{tabular}{@{}>{\raggedright\arraybackslash}p{0.2\linewidth}>{\raggedright\arraybackslash}p{0.34\linewidth}>{\raggedright\arraybackslash}p{0.17\linewidth}>{\raggedright\arraybackslash}p{0.23\linewidth}@{}}
\toprule
Result & In words & Assumptions & Status \\
\midrule
1.~The posterior finds the true model (Thm~\ref{thm:post-random}) &
The posterior mass outside the true model is of order $\rn$; the highest posterior and median probability models equal
the true model with probability tending to one. &
\as{GL}, \as{DIM}, \as{SP}, \as{SGR}, \as{FSb}, \as{EFS} &
Proved and checked (the sub-Gaussian form recorded 11-17). \\
2.~It finds the true distribution (Thm~\ref{thm:dist}) &
The posterior contracts around the true conditional law of the response, at radius of order $\rn$ in Hellinger
distance; no lower bound on the coefficients. &
Assumption~\ref{asm:dist} &
Proved and checked. \\
3.~The reported model is the true model (Thm~\ref{thm:reported}) &
For block and AR designs whose true coefficients do not cancel, with probability tending to one every run that produces
output reports exactly the true model, whatever its seed. &
\as{GR}, \as{DIM}, \as{FSt}, \as{V}, \as{EFS}, Assumption~\ref{asm:reported} &
Proved and checked; no open condition. \\
4.~Certified screening (Thm~\ref{thm:screen}) &
A run rarely certifies a reported model of at most $k_n$ predictors that misses a true predictor. &
fixed design; floor and signal bounds &
Proved and checked. \\
5.~The evidence estimate is unbiased (Prop.~\ref{prop:unbiased}) &
Given the calibration, the evidence estimate is unbiased, on every design. &
\as{C3} only &
Proved and checked. \\
6.~The rate cannot be improved (Thm~\ref{rem:sharp}) &
No design gives a posterior rate faster than $n^{-1/2}$; on two families of designs the expected posterior miss is at
least of order $\rn$. &
none for (a) &
Proved and checked. \\
\bottomrule
\end{tabular}
\end{table}

% =====================================================================================================
\subsection{Result 1: the posterior finds the true model}\label{sec:model-result}

\inwords{If the true model is small, its coefficients are fixed and nonzero, and no small group of predictors is
collinear, then the posterior puts all but a fraction of order $\rn$ of its mass on the true model. The highest posterior
model and the median probability model then equal the true model with probability tending to one, and every inclusion
probability is within order $\rn$ of $1$ for a true predictor and of $0$ for a null one.}

\begin{theorem}[Posterior concentration, random sub-Gaussian design]\label{thm:post-random}
Assume \as{GL}, \as{DIM} with $p=\praw$ (the bound on the number of raw predictors), \as{SP}, \as{SGR}, \as{FSb} and
\as{EFS}. Columns keep their raw labels, and $\Sz:=\{j\in\set{\praw}:\beta^*_j\neq0\}$ with $\bstar$ the coefficients of
the raw design; on the event that the method removes a column, set $\PIP_j:=0$ for a removed column $j$ and
$\pis{\Sz}:=0$ if $\Sz\notin\Om$, and, if no column is retained, take the posterior to be a point mass at $\emptyset$.
Then, under the joint law of the design and the noise,
\begin{gather}
  1-\pis{\Sz}=\OP\Big(\sqrt{\frac{\log n}{n}}\Big),\qquad \P\big(\HPM=\MPM=\Sz\big)\to1,\label{eq:post-random}\\
  \max_{j\in\set{\praw}}\big|\PIP_j-\ind{j\in\Sz}\big|=\OP\Big(\sqrt{\frac{\log n}{n}}\Big).\nonumber
\end{gather}
\end{theorem}

\status{Sub-Gaussian form, as stated: \textbf{proved and independently checked}, recorded on 2026-11-17, after the
freeze (an external blind audit found nothing); all its inputs are proved and independently checked, so it carries no
inherited condition. The Gaussian case $L=1$ is the recorded row 1.2: \textbf{proved and independently checked}.}

\whynat{It is the basic guarantee one expects of a Bayesian selection method with a multiplicity-correcting prior: with
enough data the posterior piles onto the true model. The size of the miss is also what one expects. It is carried by
the models that add one null predictor to the true model: each pays a prior factor of about $(\sz+1)/(2p)$ and an Occam
factor of about $n^{-1/2}$, while the best of the $p$ null predictors gains about $\log p$ in log fit; together these models
weigh of order $\sqrt{\log p/n}$ relative to the true model, and $\log p\asymp\log n$ under \as{DIM}. (This is a
heuristic; Result~6 makes it precise on two families of designs.) The last sentence of the theorem is Lemma~\ref{lem:summaries}.}

\vssl{Their Theorem~3 (p.~956) gives sensitivity (every true predictor is selected) with an exponentially small error, and
does not control specificity; ours gives exact selection, specificity included, at a polynomial rate; neither implies the
other. Their Theorem~S (supp.\ p.~13) states consistency of the posterior mode, for a different posterior, with no rate.}

\scope{Every design of our simulation studies, strong common factors included. A fixed-design version, with conditions on
the observed sample and a beta-min condition, is Theorem~\ref{thm:post-fixed} in the supplement.}

% =====================================================================================================
\subsection{Result 2: the posterior finds the true distribution}\label{sec:distribution}

\inwords{Suppose the observations are i.i.d., the response follows a sparse Gaussian linear model, the predictors have
light tails, no small group of predictors is almost exactly collinear, and no single predictor carries an unbounded
signal. Then, for every large enough level $M$, with probability at least $1-n^{-c_{\mathrm H}M}$ the posterior puts
mass at most $n^{-c_{\mathrm H}M}$ on conditional laws of the response farther than $\sqrt{M\log n/n}$ from the true
one, in Hellinger distance. No lower bound on the coefficients is needed.}

\paragraph{Setting.} As in Section~\ref{sec:a-dist}, with data set $\Dn$ and its law $\P^*$; the method runs on the
predictors that are not constant in the sample. The posterior is not defined if fewer than two predictors remain or the
response is constant in the sample; such data sets are counted in the exceptional event of the theorem. A posterior draw
$(S,\beta,\sigma^2)$, mapped back to the original units, defines the conditional law $\Normal(\bar y+(x-\bar
x)^\top\beta,\sigma^2)$ of $y$ given $x$; the law of $x$ is not estimated. Let $f$ be the joint density of $(x,y)$ with
$x$ from its true law and $y\mid x$ from this fitted law, and $f^*$ the true joint density; then
$d_H(f,f^*):=(1-\int\sqrt{ff^*})^{1/2}$ is the Hellinger distance between the fitted and the true conditional laws,
averaged over a new $x$. $\Pi(\cdot\mid\Dn)$ is the posterior of $(S,\beta,\sigma^2)$.

\begin{theorem}[Distribution-level consistency]\label{thm:dist}
There are constants $c_{\mathrm H}>0$, $M_0$ and $n_0$, depending only on $(s,A,L,c_\psi,\bar b)$, such that the following
holds. Fix a level $M\ge M_0$ and assume Assumption~\ref{asm:dist}, with the window $q_n$ of this $M$. Then for every
$n\ge n_0$ with $\sqrt n/\log n\ge M$,
\begin{equation}\label{eq:dist-i}
  \P^*\Big[\,\Pi\Big(d_H(f,f^*)>\sqrt{M\log n/n}\ \Big|\ \Dn\Big)>n^{-c_{\mathrm H} M}\Big]\ \le\ n^{-c_{\mathrm H} M}.
\end{equation}
\end{theorem}

\status{\textbf{Proved and independently checked}. Registry row 2.1.}

\whynat{It is the analogue, for this prior, of the classical posterior contraction theorems: the radius $\rn$ is the
parametric rate up to the logarithmic factor that is the usual price of not knowing the model. Prediction needs no
beta-min condition, so the posterior describes the response correctly even when some true predictors are too weak to be
selected.}

\vssl{With $\epsilon_n^2=M\log n/n$, the bound $\exp(-c\,n\epsilon_n^2)$ of their display~(15) (their Theorem~1, p.~952;
their Theorem~2 has the same form, p.~953) equals $n^{-cM}$, so~\eqref{eq:dist-i} is of the same exponential form, for
levels $M\le\sqrt n/\log n$. It uses the same distance (their supplement's convention;
their main text's distance is larger by the factor $\sqrt2$); it reaches a radius of order $\rn$, which their conditions do not
reach at polynomial $p$; larger radii and exponential $p$ are outside ours.}

\scope{Every design of our simulation studies.}

% =====================================================================================================
\subsection{Result 3: the reported model is the true model}\label{sec:reported}

\inwords{If no true predictor is masked by $15$ or fewer others, the predictors that can imitate the signal form a set
$Q$ of at most $n^b$ predictors with $b<\tfrac12$, and regressions within $Q$ are local, then there is an event,
determined by the data alone and of probability tending to one, on which every run that produces output reports exactly
the true model, whatever its seed and the outcome of its calibration.}

\begin{theorem}[The reported model equals the true model]\label{thm:reported}
Let $\Sig$ be positive definite and assume \as{GR}, \as{DIM}, \as{FSt}, \as{V}, \as{EFS}, Assumption~\ref{asm:reported}
and Convention~\as{C1}. Then, under the joint law of the design and the noise, there are data events $\mathcal G_n$ with
$\P(\mathcal G_n)\to1$ on which \emph{every} admissible run that produces output reports $\Shat=\Sz$, whatever its seed and
the outcome of its internal calibration. In particular,
\begin{equation}\label{eq:reported}
  \P(\text{the run produces output and }\Shat\neq\Sz)\to0 .
\end{equation}
If moreover Convention~\as{C2p} holds, then
\begin{equation}\label{eq:delivery}
  \P(\text{the run produces output})=1\ \ \text{for every admissible }n\ge3,\qquad \P(\Shat=\Sz)\to1 .
\end{equation}
\end{theorem}

\status{\textbf{Proved and independently checked}, with no open condition: the version proved in the registry is a
stronger form of this statement, and the two code-reading audits that its delivery route takes as hypotheses are
independently checked. That the integer limits of \as{C2p} cannot be reached on any representable input: \textbf{proved
and independently checked}. Registry row 3.1.}

\whynat{It follows the mechanism of Example~\ref{ex:leak}. The start law scores each block exactly, so it sees every
visible true predictor; its only errors are predictors that leak, all inside $Q$, so the start mode is the true model plus
some members of $Q$. Deleting a leaked predictor raises the score, so the deterministic local search from the start mode
walks down to $\Sz$, and by Result~1 nothing the run scores beats $\Sz$. This is what one would hope a well-designed
search does; the theorem says when it provably does, at the effort the software actually spends and for every seed.}

\vssl{Their statements about what their sampler reports are limits as the number of iterations grows, at fixed data and
for every data set; ours holds at the budget the software uses, jointly with $n$, under design conditions. Neither implies
the other.}

\scope{Block and AR designs whose true coefficients do not cancel (any fixed AR correlation; the constants degrade as
$\abs\rho\to1$), and independent predictors: S1--S4, S6, IND, AR and BLK; not pervasive factors with $p\gg\sqrt n$. The
supplement's Corollary~\ref{cor:famc} collects sufficient conditions for block and AR designs into one statement.}

% =====================================================================================================
\subsection{Result 4: certified screening, the guarantee that needs neither condition}\label{sec:screen}

\inwords{If small sets of observed predictors have an eigenvalue floor, the true effects are bounded by fixed constants,
and the size budget $k_n$ is small compared with $n/\log p$, then the probability that a run produces output, certifies a
reported model of at most $k_n$ predictors, \emph{and} that model misses a true predictor tends to zero. A certified model
may contain null predictors, and nothing is claimed conditionally on certification. Neither visibility nor a leakage set is
needed.}

Here \emph{certified} means that the run set its local-optimality flag (Section~\ref{sec:outputs-method}).

\begin{theorem}[Certified screening]\label{thm:screen}
Let the design be fixed and assume \as{GL} and Convention~\as{C1}; \as{SP} is not assumed. Let $k_n\in\{0,\dots,K-1\}$ be
a model-size budget, and assume $\lmin(\Rhat_S)\ge\kappa$, for a fixed $\kappa>0$, for every $S$ with $\abs{S}\le
k_n+\sz$, $\tmin\ge\tlow>0$ and $\norm{\tstar}_1\le M_1$, with $\kappa$, $\tlow$, $M_1$, $\sigs^2$ fixed, and
$(k_n+1)\log p+\log n\le c_{\mathrm{SB}}\,n$, where $c_{\mathrm{SB}}:=\kappa^2\tlow^2/\{200(M_1^2+\sigs^2)\}$. Then
\[
  \P\big(\text{the run produces output},\ \Shat\text{ is certified},\ \abs{\Shat}\le k_n,\ \Sz\not\subseteq\Shat\big)\to0 .
\]
\end{theorem}

\status{\textbf{Proved and independently checked}. Registry row 3.7.}

\whynat{A certified model is a local optimum of the posterior: no addition, deletion or swap improves it. If it missed a
true predictor, adding that predictor would bring a fit gain of order $n$, far above the prior and Occam costs of order
$\log p$; a small certified model therefore cannot miss one. It is a sure-screening property, the property a statistician
asks of a first-stage selector.}

\vssl{The nearest result is their sure-screening Theorem~3 (p.~956), for their posterior: every true predictor is
included with probability tending to one. Ours is about the model the software certifies, at its own budget.}

\begin{remark}[What is not claimed, and how often runs certify]\label{rem:screen}
(i) The bound is on an intersection; it is not a guarantee conditional on certification, no run is claimed to certify a
small model, and the certified model is not claimed to equal the true model.
(ii) \emph{In existing numerics (not a theorem),} every completed run certified its reported model: 2,394 of 2,394 in the
2,400-fit study, 376 of 376 in the strong-factor study, and 71 of 71 in a screening study. On the strong-factor designs with
equicorrelation $\varrho\in\{0.5,0.7\}$, $p\ge2.5n$ and $\sz\in\{4,8\}$, the certified model missed at least one true
predictor in 61 of 122 fits at $n=50$, 2 of 20 at $n=75$ and none of 30 at $n=100$; at the constants of these designs the
size budget of Theorem~\ref{thm:screen} requires $n\gtrsim1{,}450$.
\end{remark}

\scope{Fixed designs with the floor and signal bounds; no visibility, leakage, start or block condition.}

% =====================================================================================================
\subsection{Result 5: the evidence estimate is unbiased, on every design}\label{sec:estimates}\phantomsection\label{sec:unbiased}

\inwords{The method reports an evidence estimate, the pooled average of its four islands' estimates. Given the
calibration, it is unbiased for the evidence $Z$, on every design and data set.}

\begin{proposition}[Each island's evidence estimate is unbiased]\label{prop:unbiased}
Under Convention~\as{C3}, given the calibration, each island's evidence estimate $\Zhat_r$ is unbiased for $Z$, for every
design and data set.
\end{proposition}

By linearity the equal-weight pooled estimate~\eqref{eq:pool} is then unbiased too: on $\Acal$,
$\E[\Zhat\mid\text{calibration}]=Z$.

\status{\textbf{Proved and independently checked} (registry row 3.5), for each island and for the pooled form. It uses
two results on the coded sampler, proved under their named premises and independently checked; their final line checks
passed and were recorded on 2026-11-17, after the freeze. The frozen-run
identities behind both forms are proved and independently checked against the code of record.}

\whynat{This is the classical unbiasedness of sequential Monte Carlo estimates of normalising constants
\citep{DelMoral2004}. What is new is only that it holds for the code as written, adaptive calibration included: the
calibration is frozen before the final run draws its random numbers, so given it the final run is a fixed SMC sampler. It
is the one reported estimate with a guarantee that needs no condition on the design, and it underlies the reported lower
bound $\log(0.05\,\Zhat)$ for $\log Z$.}

\scope{Every design.}

% =====================================================================================================
\subsection{Result 6: the rate cannot be improved}\label{sec:sharp}

\inwords{(a) On every design, no model of at most $\min(K-1,p/2)$ predictors has posterior mass closer to one than about
$(\abs{S}+1)/(3\sqrt n)$, so no design gives a rate faster than $n^{-1/2}$. (b) On two families of designs the expected
posterior miss is at least of order $\rn$, the rate of Result~1.}

\begin{theorem}[The rate cannot be improved]\label{rem:sharp}
\textup{(a)} \emph{Every design.} For every $n\ge3$, every standardised design with $p\ge2$, every response and every
model $S$ with $\abs{S}\le\min(K-1,p/2)$,
\begin{equation}\label{eq:sharp-a}
  \pis{S}\ \le\ \Big(1+\frac{\abs{S}+1}{3\sqrt n}\Big)^{-1}.
\end{equation}
\textup{(b)} \emph{Witness designs.} Assume \as{GL}, that $\sz$ is fixed, the energy bounds $\norm{\mu}^2\le nM_\mu$ and
$\norm{\tstar}^2\le nM_\theta$ for fixed $M_\mu$, $M_\theta$, and that the null predictors are either orthogonal (to each
other and to the true ones) or have i.i.d.\ $\Normal(0,1)$ raw entries independent of the true predictors and of the
noise, with $p\le n^A$. Then, for every $p>\sz$, bounded or growing, with $\pi=3.14\ldots$ in this display,
\begin{equation}\label{eq:sharp-b}
  \E\big[1-\pis{\Sz}\big]\ \ge\ (1-o(1))\,(\sz+1)\,\frac{2(p-\sz)}{2p-1-\sz}\,
  \sqrt{\frac{\log\!\big(p\sqrt n/(\sz+1)\big)}{\pi n}},
\end{equation}
and, if moreover $p\to\infty$, $1-\pis{\Sz}\ge(1-\oP(1))\,(\sz+1)\sqrt{\log p/(\pi n)}$.
\end{theorem}

\status{(a) and (b): \textbf{proved and independently checked}. Registry row 1.3. On fully orthogonal designs with fixed
nonzero coefficients the supplement's expectation bound~\eqref{eq:post-exp} applies too, so there the expected miss is of
exact order $\rn$; that statement has the status of~\eqref{eq:post-exp}: proved and independently checked since 2026-11-17.}

\whynat{No fixed prior can put all its mass on one model faster than the Occam factor of one extra predictor,
$n^{-1/2}$, allows; the factor $\sqrt{\log n}$ appears because the best of $p$ null predictors gains about $\log p$ in fit.
It also explains what users see at moderate $n$: the posterior spreads over the true model and many one-predictor
extensions of it (supplement, Remark~\ref{rem:moderate-n}).}

\scope{(a) Every design; (b) the witness designs only.}
